%--------------------------------------------------------------------------------------%--------------------------------------------------------------------------------------
%
%  Global settings, dont change it! (excapt additional \usepackage commands)
%  Always use PDFLatex!
%
%--------------------------------------------------------------------------------------%--------------------------------------------------------------------------------------
\documentclass[a4paper, 12pt, oneside, BCOR1cm,toc=chapterentrywithdots]{scrbook}

\usepackage{graphicx}           % use for pdfLatex
\usepackage{makeidx} % f\"{u}r Benutzung des Befehls \printindex
\usepackage[colorlinks=false]{hyperref}
\usepackage{tocbibind}
\usepackage{blindtext}
\usepackage{subfigure} 
\usepackage{acronym}
\usepackage{minted}
\usepackage{multicol}
\usepackage[dvipsnames]{xcolor}
\usepackage{enumitem}
\usepackage{tcolorbox}
\usepackage{listings}

\newenvironment{myreq}[1]{%
\setlist[description]{font=\normalfont\color{darkgray}}%
\begin{tcolorbox}[colframe=black,colback=white, sharp corners, boxrule=1pt]%
\bfseries\color{blue}%
\begin{description}#1}%
{\end{description}\end{tcolorbox}}

\newcommand{\threeinline}[3]{\begin{multicols}{3}#1 #2 #3\end{multicols}}
\newcommand{\twoinline}[2]{\begin{multicols}{2}#1 #2\end{multicols}}

\newcommand{\reqno}{\item[Requirement-ID:]}
\newcommand{\reqtype}{\item[Requirement Type:]}
\newcommand{\reqevent}{\item[Event:]}
\newcommand{\reqdesc}{\item[Description:]}
\newcommand{\reqrat}{\item[Rationale:]}
\newcommand{\reqorig}{\item[Originator:]}
\newcommand{\reqfit}{\item[Fit Criterion:]}
\newcommand{\reqsatis}{\item[Customer Satisfaction:]}
\newcommand{\reqdissat}{\item[Customer Dissatisfaction:]}
\newcommand{\reqdep}{\item[Dependencies:]}
\newcommand{\reqconf}{\item[Conflicts:]}
\newcommand{\reqmater}{\item[Materials:]}
\newcommand{\reqhist}{\item[History:]}




\hypersetup{%
bookmarksnumbered=true, hyperindex=true,
%
%Im Acrobat Reader Subtitel 1. Ebene anzeigen
bookmarksopen=true, bookmarksopenlevel=1,
%
pdfborder=0 0 0 % Keine Box um die Links!
}

% --------------------------------------------------------------
% Force Tables and List to be added in Table of Content
% --------------------------------------------------------------

\renewcommand*{\tableofcontents}{%
  	\begingroup
  	\tocsection
  	\tocfile{\contentsname}{toc}
  	\endgroup
}
\renewcommand*{\listoffigures}{%
  	\begingroup
  	\tocsection
  	\tocfile{\listfigurename}{lof}
  	\endgroup
}
\renewcommand{\listoftables}{
	\begingroup
	\tocsection
	\tocfile{\listtablename}{lot}
	\endgroup
}
\begin{document}

%--------------------------------------------------------------------------------------%--------------------------------------------------------------------------------------
%
%  Here starts the userspace !
%
%--------------------------------------------------------------------------------------%--------------------------------------------------------------------------------------

%--------titlepage
\begin{titlepage}

{
    \begin{center}
        \raisebox{-1ex}{\includegraphics[scale=1.5]{TU_Chemnitz_positiv_gruen.pdf}}\\
    \end{center}
    \vspace{0.5cm}
}

\begin{center}

\LARGE{\textbf{Titel des Abschlussberichts}}\\
\vspace{1cm}


\Large{\textbf{Softwaretechnikpraktikum - Abschlussbericht}}\\ 
\vspace{1cm}

Fakultät für Informatik\\
Professur Softwaretechnik
\end{center}
\vspace{3cm}
Eingereicht von: Gruppe X (Mitglied 1, Mitglied 2...)\\
Einreichungsdatum: 12.12.1212\\
\vspace{0.3cm}\\
Betreuerin: Prof. Dr. Janet Siegmund \\
Betreuer: Dominik Gorgosch

\end{titlepage}

%---------------------------------------------------------
% Table of Contents, List of figures, List of Tables
%---------------------------------------------------------

\renewcommand{\contentsname}{Inhaltsverzeichnis}
\renewcommand{\listfigurename}{Abbildungsverzeichnis}
\renewcommand{\listtablename}{Tabellenverzeichnis}
\renewcommand{\figurename}{Abbildung}
\tableofcontents
\listoffigures
\listoftables

%---------------------------------------------------------
% Abkürzungsverzeichnis
%---------------------------------------------------------
\twocolumn
\addchap{Abkürzungsverzeichnis}
\begin{acronym}[Bash]
 \acro{TUC}{TU Chemnitz}

\end{acronym}

\onecolumn
%---------------------------------------------------------
% Here starts the real work
%---------------------------------------------------------

\chapter{Einführung}
% example of acronym
Hier sehen Sie, wie Sie eine Abkürzung nutzen können. Dieses Modul findet an der \ac{TUC} statt.
  % Load Data from File intro.tex

\chapter{Tabellen}
% Example of, how to use a Table

In diesem Kapitel sehen sie ein Beispiel für eine Tabelle. Es gibt deutlich komplexere Tabellen in Latex und auch weitere Packages dafür. Wenn Sie Fragen haben, melden Sie sich gerne bei mir.

\begin{center}
\begin{tabular}{|c|c|c|c|}
\hline
Wert 1 & Wert 2 & Wert 3 & Wert 4\\
in Einheit1 & in Einheit2 & in Einheit3 & in Einheit4 \\
\hline
192 & 80& 0.3153 & 0.4900\\
500 & 120& 0.1229& 0.1787\\
1000 & 120& 0.0680& 0.0880\\
2000 & 120& 0.0361& 0.0441\\
5000 & 140& 0.0256& 0.0305\\
5000 & 164& 0.0343& 0.0880\\
\hline
\end{tabular}
\captionof{table}{This is the caption of the table}
\label{tab:table1}
\end{center}

 % Load Data from File example_tables

\chapter{Abbildungen}
% Example of, how to use figures

Hier sehen Sie die Einbindung von Abbildungen/Grafiken/Bildern.

\begin{figure}[h]
\centering
\includegraphics[width=0.6\textwidth]{TU_Chemnitz_positiv_gruen.pdf}
\caption{Abbildung 1}
\label{fig:pic0}
\end{figure}

Wenn Sie mehrere Bilder nebeneinander einfügen wollen, können Sie eine subfigure oder minipage nutzen.

\begin{figure}[h]
    \subfigure[Dies ist die erste Grafik]{\includegraphics[width=0.49\textwidth]{TU_Chemnitz_positiv_gruen.pdf} \label{fig:pic1}}
    \subfigure[Dies ist die zweite Grafik]{\includegraphics[width=0.49\textwidth]{TU_Chemnitz_positiv_gruen.pdf}\label{fig:pic2}}
\caption{Dies ist die Beschriftung der gesamten Grafik}
\end{figure}

Für die Referenzierung eines Bildes, nutzen Sie folgenden Befehl~\ref{fig:pic1}.
 % Load Data from File example_figures

\chapter{Codebeispiel}
%how to implement code
%https://de.overleaf.com/learn/latex/Code_Highlighting_with_minted

Im folgenden Kapitel sehen Sie, wie Sie Code in Latex einbinden können. Eine weitere Option wäre das Package listings. Zuerst sehen Sie ein Python-Beispiel.

\begin{minted}{python}
def addNum(list):
    sum = 0
    
    for val in list:
        sum = sum+val
        
    return sum
\end{minted}

Hier sehen Sie ein Java-Beispiel:

\begin{minted}{java}
public class Calculator {
    public static int addNum(int[] list) {
          int sum = 0;

        for(int i = 0; i < list.length; i++) {
            sum = sum + list[i];
        }

        return sum;
    }

    public static void main(String[] args) {
        int[] numbers = {1,2,3,4,5,6,7,8,9};
        int sum = addNum(numbers);
        System.out.printf("Sum: %d", sum);
    }
}
\end{minted}


Wenn Sie ihren Code referenzieren möchten, dann müssen Sie Listings nutzen. Dies geht auch zusammen mit minted. 

\begin{listing}
\begin{minted}[mathescape,
               linenos,
               numbersep=5pt,
               gobble=2,
               frame=lines,
               framesep=2mm]{java}
    public class Example {
        public static void main(String[] args) {
            int[] array = {1, 5, 3, 2, 4};
            
            for (int i = 0; i < array.length; i++) {
                System.out.println(array[i]);
            }
        }
    }
\end{minted}
\caption{Array in Java}
\label{java_example}
\end{listing}
\newpage

In line~2 of listing \ref{java_example}, we can see that...

%import code from file:
%\inputminted{python}{src/example.py} % Load Data from File example_code

\chapter{Abschnitte und Zitieren} % Load Data from File example_sections
\section{Abschnitt 1}
Nach McCracken et al. bestehen folgende Probleme ... bei Programmieranfängern~\cite{McCracken2001}.

\section{Abschnitt 2}
\subsection{Unterabschnitt 2.1}
\subsection{Unterabschnitt 2.2}



\chapter{Volere} % Load Data from File volere_template
%Quelle: https://tex.stackexchange.com/questions/515525/how-to-make-a-custom-requirement-card-like-voleres

\begin{myreq}
  \threeinline
    {\reqno 75}
    {\reqtype Funktional}
    {\reqevent -}
  \reqdesc description text description text description text 
  \reqrat some more text some more text some more text 
  \reqorig other text other text other text 
  \reqfit longer text that needs more than one line longer text that needs more than one line
  \twoinline
    {\reqsatis 5}
    {\reqdissat 3}
  \twoinline
  {\reqdep some more text}
  {\reqconf 111}
  \reqmater other text other text other text 
  \reqhist other text other text other text 
\end{myreq}


%---------------------------------------------------------
% bibliography based on Springer Design
%---------------------------------------------------------

\bibliographystyle{splncs03}
\renewcommand{\bibname}{Literaturverzeichnis}
\bibliography{bibliography}

\printindex

\end{document}
